\documentclass[11pt]{article}
\usepackage[margin=1in]{geometry}
\usepackage{amsmath,amssymb}
\usepackage[hidelinks]{hyperref}

\begin{document}

\section*{Human Review: No sparse Euclidean grid in \texorpdfstring{$\ell_\infty^n$}{ell-infinity-n}}

\noindent Original automatic proof: \texttt{solution\_packet.pdf}

\noindent Checked date: 18 June 2026

\noindent Status: Manually checked.

\noindent Verdict: The proof appears correct.

\section*{Human Comments}

The proposed argument appears to give a correct negative answer to Fresen's
Problem 4, already for the spaces \(X=\ell_\infty^n\).  The problem asks
whether one can choose bases in all \(n\)-dimensional Banach spaces so that
sufficiently sparse coefficient vectors are uniformly Euclidean.  The proof
shows that this is impossible in \(\ell_\infty^n\) for any fixed distortion
constant \(C\), even if the allowed sparsity tends to infinity arbitrarily
slowly.

The translation to matrices is clear.  A basis of \(\ell_\infty^n\) gives a
linear map
\[
  T:\mathbb R^n\longrightarrow \ell_\infty^n
\]
with rows \(r_i=(r_{ij})_{j=1}^n\).  The upper sparse estimate implies that,
for every row \(i\) and every coordinate set \(A\) with \(|A|\leq k\),
\[
  \left(\sum_{j\in A} r_{ij}^2\right)^{1/2}\leq C.
\]
This follows by viewing the \(i\)-th row as a functional on the Euclidean
coordinate subspace supported on \(A\).

\section*{Counting Lemma}

The counting lemma is the main combinatorial ingredient and appears correct.
It says that \(2n\) subsets of the signed coordinate set
\[
  \Omega_n=\{1,\ldots,n\}\times\{-1,1\},
\]
each of size at most \(Mh\), cannot substantially hit every signed \(h\)-set
once \(h\to\infty\) and \(h\leq \sqrt{\log n}\).

For a fixed set \(H\subseteq \Omega_n\) with \(|H|\leq Mh\), the number of
signed \(h\)-sets meeting \(H\) in at least \(q_0=\lceil \eta h\rceil\)
points is bounded by
\[
  \sum_{q=q_0}^h
  \binom{Mh}{q}\binom{n}{h-q}2^{h-q}.
\]
Dividing by the total number \(\binom nh2^h\) of signed \(h\)-sets gives the
fraction hit by this particular \(H\).  Using
\[
  \frac{\binom{n}{h-q}}{\binom nh}
  \leq
  \left(\frac{2h}{n}\right)^q
\]
for \(h\leq n/2\), this fraction is at most a constant-factor variant of
\[
  h\,2^{Mh}\left(\frac{h}{n}\right)^{\eta h}.
\]
After multiplying by \(2n\), the logarithm is bounded above by
\[
  \log(2n)+\log h+Mh\log 2
  -
  \eta h\log\left(\frac nh\right),
\]
which tends to \(-\infty\).  The negative term is essentially
\(-\eta h\log n\), and it dominates the positive terms because
\(h\to\infty\) while \(h\leq \sqrt{\log n}\).  Hence the union of the \(2n\)
families cannot cover all signed \(h\)-sets.

This part of the proof is correct, although it would be useful in a polished
version to expand the counting estimate slightly more than in the current
automatic output.

\section*{Application to the Counterexample}

The proof then chooses
\[
  h=
  \left\lfloor
  \min\left\{\frac{k}{64C^2},\sqrt{\log n}\right\}
  \right\rfloor,
  \qquad
  \theta=\frac{1}{4\sqrt h}.
\]
Since \(k\to\infty\) and \(n\to\infty\), one has \(h\to\infty\) directly;
there is no real need to pass to a subsequence, apart from discarding
finitely many small values if desired.

For each row \(i\), set
\[
  A_i=\{j: |r_{ij}|\geq \theta\}.
\]
The proof first shows that \(|A_i|<k\), because otherwise one could choose
\(k\) elements of \(A_i\) and contradict the row mass estimate.  This step is
needed because the row mass estimate is only known for coordinate sets of
size at most \(k\).  Once \(|A_i|<k\) is known, the same estimate applied to
\(A_i\) gives
\[
  |A_i|\leq 16C^2h.
\]

For each row \(i\) and sign \(s\in\{-1,1\}\), define
\[
  H_i^s
  =
  \{ (j,\varepsilon)\in\Omega_n:
      s\varepsilon r_{ij}\geq \theta \}.
\]
Then \(|H_i^s|\leq 16C^2h\), so these are precisely the small signed patterns
to which the counting lemma will be applied.

Given an arbitrary signed \(h\)-set
\[
  E(S,\varepsilon)=\{(j,\varepsilon_j):j\in S\},
\]
consider the vector \(a_j=\varepsilon_j/\sqrt h\) on \(S\) and \(a_j=0\)
outside \(S\).  The lower estimate gives a row \(i\) and a sign
\(s\in\{-1,1\}\) such that
\[
  \sum_{j\in S}s r_{ij}\varepsilon_j\geq \sqrt h.
\]
Now define
\[
  B=\{j\in S: s r_{ij}\varepsilon_j\geq \theta\}.
\]
It is helpful to record explicitly that
\[
  E(S,\varepsilon)\cap H_i^s
  =
  \{(j,\varepsilon_j):j\in B\}.
\]
Thus showing that \(|B|\) is large is exactly the same as showing that the
signed \(h\)-set \(E(S,\varepsilon)\) has large intersection with one of the
sets \(H_i^s\).

The terms outside \(B\) contribute at most \(h\theta=\sqrt h/4\), so the
contribution from \(B\) is at least \(3\sqrt h/4\).  By Cauchy-Schwarz and
the row mass estimate on \(B\),
\[
  \frac{3\sqrt h}{4}
  \leq
  C |B|^{1/2}.
\]
Hence
\[
  |B|\geq \frac{9}{16C^2}h.
\]
Therefore every signed \(h\)-set has intersection at least
\(\eta h\), with \(\eta=9/(16C^2)\), with one of the \(2n\) sets \(H_i^s\),
each of size at most \(Mh\), with \(M=16C^2\).  This contradicts the counting
lemma.

\section*{Literature Check}

A preliminary citation-oriented search found no indexed citing papers for the
published Studia Mathematica version of Fresen's article.  OpenAlex, Crossref,
and OpenCitations all report no citing works for the journal version, and
targeted searches for the distinctive Problem 4 terminology did not reveal an
existing answer.  This does not replace author feedback or a final literature
check, but it supports treating the packet as a potentially new negative
answer.

\section*{Review Outcome}

Archive this packet as an apparently correct full negative answer to Fresen's
Problem 4, pending author feedback.  The argument is sound, but a polished
version should expand the counting lemma and make the relation
\[
  E(S,\varepsilon)\cap H_i^s
  =
  \{(j,\varepsilon_j):j\in B\}
\]
explicit when the set \(B\) is introduced.

\end{document}
